\documentclass[11pt]{article}
\usepackage[margin=1in]{geometry}
\usepackage[T1]{fontenc}
\usepackage{lmodern}
\usepackage{amsmath,amssymb}
\usepackage{hyperref}
\emergencystretch=2em

\title{Literature Status for arXiv:1208.2384: Adversarial Gowers Games}
\author{Run fa\_banach\_001, agent\_lane\_04}
\date{2026-07-18}

\begin{document}
\maketitle

\section*{Status}

This is a \texttt{literature\_already\_answered} packet.  It records that
Problem~1.2 in Christian Rosendal's \emph{Determinacy of adversarial Gowers
games}, arXiv:1208.2384, is explicitly answered by Noe de Rancourt's
\emph{Ramsey theory without pigeonhole principle and the adversarial Ramsey
principle}, arXiv:1805.04954.

\section*{Source Question}

Rosendal's Theorem~1.1 proves the adversarial Gowers-game dichotomy for
\(F_{\sigma\delta}\) and \(G_{\delta\sigma}\) subsets of \(E^\omega\).  On PDF
page~3, immediately after Theorem~1.1, Problem~1.2 asks whether
Theorem~1.1 holds for all Borel sets \(\mathcal A\), or even for analytic sets
in the presence of large cardinals.

\section*{Supporting Answer}

The supporting paper explicitly identifies the same questions.  On PDF page~6,
de Rancourt states Rosendal's Question~1.10, asking whether every Borel set is
adversarially Ramsey, and Question~1.11, asking whether every analytic set is
adversarially Ramsey in the presence of large cardinals.  The same paragraph
says that the paper proves the answer to Question~1.10 is yes, and that the
answer to Question~1.11 is yes in the presence of a measurable cardinal.

The Borel answer is Theorem~2.4, the abstract adversarial Ramsey principle:
every Borel subset of \(X^\omega\) is adversarially Ramsey.  Immediately after
the theorem, the paper says that, in the Rosendal space, the adversarial Gowers
games defined there are exactly the games from the introduction, so
Theorem~2.4 gives a positive answer to Question~1.10.

For the analytic large-cardinal clause, Corollary~2.6 says that if every
\(\Gamma\)-subset of \(\mathbb R^\omega\) is determined, then every
\(\Gamma\)-subset of \(X^\omega\) in an analytic Gowers space is adversarially
Ramsey.  The following paragraph invokes Martin's measurable-cardinal theorem
for analytic determinacy and concludes that every analytic set in an analytic
Gowers space is adversarially Ramsey; it explicitly says this answers
Question~1.11.

\section*{Scope}

This is not a new proof packet.  It records a known positive literature answer
to Rosendal's Problem~1.2.  The Borel part is answered in ZFC through Borel
determinacy.  The analytic part is answered under the measurable-cardinal
hypothesis used by de Rancourt, matching Rosendal's large-cardinal formulation
but not giving a ZFC analytic result.

The supporting paper uses its own convention for naming the adversarial games
in the introductory reformulation.  For that reason, the identification here is
based on de Rancourt's explicit Rosendal Questions~1.10 and~1.11 and his stated
positive answers, rather than on an isolated comparison of the letters \(A\)
and \(B\).

\section*{Search Evidence}

The lane-1 queue selected arXiv:1208.2384.  Cheap run indexes were searched for
the arXiv id, source title, ``adversarial Gowers games'', ``strategic Ramsey'',
``block sequences'', and ``Banach game''; no duplicate packet was found.
Targeted web search for the exact source title and for ``adversarial Gowers
games Borel analytic large cardinals'' identified arXiv:1805.04954.  The local
source TeX and extracted PDF text verify the exact problem statement, the
supporting paper's explicit Rosendal-question restatement, Theorem~2.4, and the
measurable-cardinal paragraph answering the analytic clause.

{\raggedright
\begin{thebibliography}{9}

\bibitem{Rosendal2014}
C.~Rosendal,
\emph{Determinacy of adversarial Gowers games},
arXiv:1208.2384; Fund. Math. \textbf{227}, 163--178 (2014).

\bibitem{deRancourt2020}
N.~de Rancourt,
\emph{Ramsey theory without pigeonhole principle and the adversarial Ramsey
principle},
arXiv:1805.04954; Trans. Amer. Math. Soc. \textbf{373}, 5025--5056 (2020).

\end{thebibliography}
}

\end{document}
